% Official CC BY4.0 published HTML S3; original math alttext and author text, literal original numbering.


\par\medskip\noindent\textbf{3 Outline of the proof of Theorem 1.1}\par




The proof of Theorem 1.1 is built on an efficient construction of the image of metapolynomials under projections onto the isotypic spaces.





Roughly speaking, we will construct the projectors of Theorem 1.1 as linear combinations of repeated derivations applied to metapolynomials. These derivations live naturally in a vector space called the Lie algebra $\mathfrak{gl}_{k}$ of $\textnormal{GL}_{k}$. They are defined as the linearization of linear change of coordinates by elements of $\textnormal{GL}_{k}$ applied to polynomials on which a metapolynomial is evaluated. Repeated applications of these operations is formally described using products in an associative algebra $\mathcal{U}(\mathfrak{gl}_{k})\supseteq\mathfrak{gl}_{k}$ called the universal enveloping algebra of $\mathfrak{gl}_{k}$.
More concretely, a standard basis element in $\mathfrak{gl}_{k}$ corresponds to an analogue of a shifted partial derivative for metapolynomials (see ˜4.2), and an element of $\mathcal{U}(\mathfrak{gl}_{k})$ captures repeated applications of such shifted partial derivatives.





Lie algebras and their universal enveloping algebras have a long history in mathematics, and their theory is rich and well-developed. We review the theory in Section 4 to the extent required to construct the projectors. The construction is completely explicit. The universal enveloping algebra $\mathcal{U}(\mathfrak{gl}_{k})$ contains elements $\{C_{1},\ldots,C_{r}\}$ which separate isotypic components in the following sense: each $C_{i}$ acts by scalar multiplication on each isotypic component, and for any two distinct isotypic components there is a $C_{i}$ for which the corresponding scalars are different; the precise statement is the content of Theorem 4.11 and Theorem 4.13. This property is used to construct the projectors via an analogue of standard multivariate polynomial interpolation, see Section 5.1. The elements $\{C_{1},\ldots,C_{r}\}$ are called Casimir operators, and they are defined as generators of certain commutative subalgebras of $\mathcal{U}(\mathfrak{gl}_{k})$. For each case of Theorem 1.1, there is a suitable algebra for which the corresponding isotypic components are the spaces of interest, see Table 1.





The proof that the projectors can be implemented efficiently is built on the following insights. The Poincaré–Birkhoff–Witt theorem (Theorem 4.9) allows one to express each projector as a linear combination in $\mathcal{U}(\mathfrak{gl}_{k})$ of $m$-fold products of elements of $\mathfrak{gl}_{k}$ with $m$ bounded polynomially in the number of isotypic components. This number is exponential in $k$, which is polylogarithmic in $N$ as observed in Section 2.5. We use this fact to construct circuits of quasipolynomial size for the projectors in Section 5.2. In turn, this will complete the proof of Theorem 1.1.



